% Family 5 visual chapter. Included by the monograph and the visual guide.
% Exact proof sources: research_v11/coupled_growth_proof.md,
% coupled_identifiability.md, and coupled_structure.md.

\pageheading{FAMILY 5}{Two functions exchange growth}
For four real contractions $0<q,r,a,b<1$, define the normalized pair
\[
\boxed{\begin{aligned}
U'(z)&=1+U(qz)V(rz),\\
V'(z)&=1+U(az)V(bz),\qquad U(0)=V(0)=0.
\end{aligned}}
\]
Each equation combines one value of $U$ with one value of $V$. The four
scales specify where that exchange occurs. Both functions begin with the
same unit slope, while later coefficients retain the two equations' different
scale addresses.

\vfig[0.96]{coupled_01_matrix}{The worked pair
$(q,r,a,b)=(1/5,2/5,3/5,7/10)$ has exact order one. The two displayed
identities give the positive growth vector $(1/2,1)^T$.}

There is a unique entire solution pair. Both functions are odd, and every
odd Taylor coefficient is positive. The common order is governed by the
positive matrix
\[
P(s)=\begin{pmatrix}q^s&r^s\\a^s&b^s\end{pmatrix},
\qquad \operatorname{spr}P(\rho)=1.
\]
The eigenvector supplies the relative weights that a row sum misses. In
the example, the row sums at $s=1$ are $3/5$ and $13/10$. The coupled
system nevertheless has order exactly one in both components.

\pageheading{FAMILY 5 / CONSTRUCTION}{The coefficients preserve both addresses}
Write $U(z)=\sum_{n\ge0}u_nz^{2n+1}$ and
$V(z)=\sum_{n\ge0}v_nz^{2n+1}$, with $u_0=v_0=1$.
The first equation gives the exact recursion
\[
u_n=\frac1{2n+1}\sum_{i+j=n-1}
q^{2i+1}r^{2j+1}u_i v_j,\qquad n\ge1.
\]
The second replaces $(q,r)$ by $(a,b)$ while retaining the same products
$u_i v_j$. Thus each new coefficient uses the shared history of the pair.

\vfig[0.97]{coupled_02_functions}{For
$(q,r,a,b)=(1/3,2/3,1/2,3/4)$, the components separate despite their
common initial slope. These numerical curves use odd indices through 240.
The two truncations at indices 120 and 240 agree at $x=8$ to displayed
double precision.}

The first distinction already appears at degree three:
\[
U(z)=z+\frac{qr}{3}z^3+O(z^5),\qquad
V(z)=z+\frac{ab}{3}z^3+O(z^5).
\]
Induction gives $0<u_n,v_n\le2^{-n}$, which supplies a convergent germ.
If both functions are known on a disk of radius $R$, their integral
equations extend them to radius $R/Q$, where $Q=\max(q,r,a,b)<1$.
Repeating this step reaches every finite disk.

The numerical agreement above checks the plotted approximation. The later
certificate page supplies an exact bound on the omitted tail.

\pageheading{FAMILY 5 / GROWTH}{One order, two exact type proportions}
The entire functions have the same order $\rho$, with positive lower type
and finite upper type. Positivity places the maximum modulus on the positive
real axis, so $M_U(R)=U(R)$ and $M_V(R)=V(R)$.

\vfig[0.96]{coupled_03_growth}{The order-one example has different component
amplitudes. The right panel shows finite-radius slopes near one. The curves
use odd indices through 800. Agreement with index 400 is a numerical
diagnostic, and the order itself follows from the analytic proof.}

For $F_1=U$ and $F_2=V$, define
\[
\underline\tau_i=\liminf_{R\to\infty}\frac{\log M_{F_i}(R)}{R^\rho},
\qquad
\overline\tau_i=\limsup_{R\to\infty}\frac{\log M_{F_i}(R)}{R^\rho}.
\]
If $P(\rho)w=w$ with $w>0$, then both type vectors have precisely these
Perron proportions:
\[
\boxed{\underline{\boldsymbol\tau}=\underline\tau\,w,
\qquad\overline{\boldsymbol\tau}=\overline\tau\,w,
\qquad0<\underline\tau\le\overline\tau<\infty.}
\]
For the worked pair, both the lower type and the upper type of $V$ are
twice the corresponding type of $U$. Either both ordinary type limits
exist or neither does. Equality of the lower and upper types remains open.
The finite-radius slope picture does not decide that question.

\pageheading{FAMILY 5 / PARAMETER SPACE}{Growth follows a spectral surface}
The order is the unique positive solution of
\[
\boxed{(1-q^\rho)(1-b^\rho)=(ar)^\rho.}
\]
The left side increases from zero to one as $\rho$ increases. The right
side decreases from one to zero. Their unique crossing determines the
growth exponent without fitting a truncated Taylor curve.

\vfig[0.80]{coupled_04_landscape}{Numerical solutions of the proved
characteristic equation, with $r=2/5$ and $a=3/5$. The highlighted pair
$(q,b)=(1/5,7/10)$ has order exactly one by rational arithmetic.}

The order depends on the self scales $q,b$ and on the product $ar$ of
the cross scales. Holding $q,b,ar$ fixed preserves the order, while
changing the cross scales separately generally changes the functions and
their type proportions:
\[
\frac{w_1}{w_2}=\frac{r^\rho}{1-q^\rho}
=\frac{1-b^\rho}{a^\rho}.
\]
Consequently, equal order does not determine the family member. It fixes
one growth exponent within a larger parameter space. The coefficients
carry enough additional information to recover all four ordered scales
away from the reduction surfaces described below.

\pageheading{FAMILY 5 / PRECISION}{A coefficient list with a certified tail}
A finite Taylor list becomes an exact approximation certificate when it
is paired with a bound on a larger disk. Let $F$ be either component,
and let $F_N$ retain the odd indices $0,\ldots,N$.

\vfig[0.95]{coupled_05_certification}{Each curve is computed from an exact
rational upper bound for the same worked pair. Increasing the coefficient
count certifies a smaller uniform error on the stated complex disk.
The logarithmic plotting coordinates are rounded.}

Put $Q=\max(q,r,a,b)$. One elementary majorant starts with
$B(x)=x/(1-x^2)$ for $0\le x<1$ and works outward through the finite
contraction ladder using
\[
B(x)=x\bigl(1+B(Qx)^2\bigr).
\]
The integral equations show that this bounds both components on the disk
$|z|\le x$. For $S>R\ge0$, positivity then gives
\[
\boxed{\max_{|z|\le R}|F(z)-F_N(z)|
\le B(S)\left(\frac RS\right)^{2N+3}.}
\]
The exponent is the first omitted odd degree. Every later term has an
even smaller ratio between its values at radii $R$ and $S$.
For the worked pair, $N=30$, $R=1$ and $S=2$ give an error below
$1.600\times10^{-13}$ for both functions. The laboratory returns
rational interval endpoints, so the certificate survives decimal rounding.

\pageheading{FAMILY 5 / SEQUENCES}{The phase calculus transports the pair}
Use the derivative normalization
$A_n=(2n+1)!u_n$, $B_n=(2n+1)!v_n$. Its recurrence has integer
binomial weights:
\[
A_n=\sum_{i+j=n-1}\binom{2n}{2i+1}
q^{2i+1}r^{2j+1}A_iB_j,\qquad n\ge1.
\]
The second component again replaces $(q,r)$ by $(a,b)$.

\vfig[0.91]{coupled_06_polyphase}{The full derivative array of $V$ is split
into three residue lanes. Applying taps $(1,-2,1)$ with the correct block
carries, then interleaving, agrees exactly with direct convolution.}

These derivative arrays are polynomials with nonnegative integer
coefficients in the four scale variables. Their total degree is at most
$n(n+1)$. If every parameter denominator divides $D_0$, then
\[
D_0^{n(n+1)}A_n,\qquad D_0^{n(n+1)}B_n
\quad\hbox{are integers}.
\]
This uniform clearing factor makes exact rational families available to
integer sequence experiments. It need not be the smallest possible factor.

Ordinary taps on a derivative array correspond to repeated integration
shifts of its exponential generating function. Multiplication of two
exponential generating functions instead requires binomial convolution.
The catalog preserves that distinction and marks generic integer and
signed scale examples as formal series outside the proved analytic domain.

\pageheading{FAMILY 5 / REDUCTION}{Two surfaces where one component rescales to the other}
Set $p=qr$, $k=ab$, and use the ordinary Taylor coefficients $u_2,v_2$
at degree five. The exact separation quantity factors as
\[
\boxed{D=k^2u_2-p^2v_2
=\frac{pk ar}{15}(q^2-b^2)(qb-ar).}
\]

\vfig[0.88]{coupled_07_reduction_geometry}{The two reduction surfaces are
$q=b$ and $qb=ar$. Their union is exactly the zero set of $D$ in the
positive contraction domain. Outside this union, the components cannot
be obtained from each other by amplitude and argument rescaling.}

There exist nonzero complex constants $A,B$ with $V(z)=A\,U(Bz)$
if and only if $D=0$. On either surface, put $s=\sqrt{k/p}>0$.
The scalar family is defined by
$T_{x,y}'(z)=1+T_{x,y}(xz)T_{x,y}(yz)$, $T_{x,y}(0)=0$.
It gives the explicit reduction
\[
U(z)=\sqrt{s}\,T_{q,sr}(z/\sqrt{s}),\qquad
V(z)=s^{-1/2}T_{q,sr}(\sqrt{s}\,z).
\]
Both effective scales lie in $(0,1)$. Equality $U=V$ is more restrictive:
it occurs exactly when $\{q,r\}=\{a,b\}$.

This classifies relations between the normalized components. Separately,
no invertible constant linear coordinate change decouples the defining
operator into two autonomous scalar equations on every analytic input.
That operator statement does not exclude nonlinear changes of variables
or other special-function identities.

\pageheading{FAMILY 5 / RECOVERY}{Degree seven identifies the four scales}
Away from $D=0$, the six coefficients
$(u_1,v_1,u_2,v_2,u_3,v_3)$ uniquely recover the ordered parameters
$(q,r,a,b)$ in $(0,1)^4$. These are the coefficients at degrees three,
five and seven of the two labeled functions.

\vfig[0.92]{coupled_08_inverse}{Cubic terms recover the two scale products.
Fifth-degree terms leave at most two candidates for each remaining scale.
The seventh-degree terms select the unique admissible candidates when
$D\ne0$. Repeated quadratic roots are already unique.}

First $p=3u_1$ and $k=3v_1$. The possible values $x=q^2$ and $y=a^2$
satisfy
\[
p x^2-\frac{15u_2}{p}x+kp^2=0,\qquad
p y^2-\frac{15v_2}{k}y+k^3=0.
\]
For two distinct positive candidates $x,x'$, the predicted seventh
coefficients differ by
\[
u_3(x)-u_3(x')=\frac{pD}{7k^2}(x^2-x'^2)\ne0.
\]
The second component supplies the analogous separation for $y$.
Finally take the positive square roots and recover $r=p/q$, $b=k/a$.

This is an exact-data inverse theorem. Near a reduction surface, small
$D$ can make candidate predictions close, so a uniform stable inverse
does not follow. The degree-seven recovery and the reduction classification
are proved statements. Historical priority of these particular results
is addressed separately by the bounded literature comparison.
